\documentclass[11pt]{article}

\usepackage[letterpaper,margin=1in]{geometry}
\usepackage[T1]{fontenc}
\usepackage[utf8]{inputenc}
\usepackage{lmodern}
\usepackage{microtype}
\usepackage{amsmath,amssymb,mathtools,bm}
\usepackage{enumitem}
\usepackage{xcolor}
\usepackage{booktabs}
\usepackage{array}
\usepackage{hyperref}
\usepackage[most]{tcolorbox}
\usepackage{fancyhdr}
\usepackage{titlesec}
\usepackage{parskip}

\definecolor{qftblue}{HTML}{245A8D}
\definecolor{qftgreen}{HTML}{2F6B4F}
\definecolor{qftred}{HTML}{A23B3B}
\definecolor{qftgray}{HTML}{F4F6F8}
\definecolor{qftgold}{HTML}{C08B2C}

\hypersetup{
  colorlinks=true,
  linkcolor=qftblue,
  urlcolor=qftblue,
  citecolor=qftgreen,
  pdftitle={PHYS 662 Lecture 3: Quantum Sound and Phonons},
  pdfauthor={Nima Lashkari}
}

\pagestyle{fancy}
\fancyhf{}
\lhead{PHYS 662: Quantum Field Theory I}
\rhead{Lecture 3}
\cfoot{\thepage}
\setlength{\headheight}{14pt}

\titleformat{\section}{\Large\bfseries\color{qftblue}}{\thesection.}{0.6em}{}
\titleformat{\subsection}{\large\bfseries\color{qftgreen}}{\thesubsection}{0.6em}{}
\titleformat{\subsubsection}{\normalsize\bfseries\color{qftred}}{\thesubsubsection}{0.6em}{}

\setlist[itemize]{leftmargin=2em,itemsep=0.25em,topsep=0.4em}
\setlist[enumerate]{leftmargin=2em,itemsep=0.25em,topsep=0.4em}

\newcommand{\dd}{\mathrm{d}}
\newcommand{\e}{\mathrm{e}}
\newcommand{\ii}{\mathrm{i}}
\newcommand{\cH}{\mathcal{H}}
\newcommand{\cF}{\mathcal{F}}
\newcommand{\R}{\mathbb{R}}
\newcommand{\comm}[2]{\left[#1,#2\right]}
\newcommand{\pb}[2]{\left\{#1,#2\right\}_{\mathrm{PB}}}
\newcommand{\ket}[1]{\left|#1\right\rangle}
\newcommand{\bra}[1]{\left\langle#1\right|}
\newcommand{\braket}[2]{\left\langle#1\middle|#2\right\rangle}

\newtcolorbox{keyidea}[1][]{
  enhanced,
  breakable,
  colback=qftgray,
  colframe=qftblue,
  boxrule=0.8pt,
  arc=1.5mm,
  title={#1},
  fonttitle=\bfseries,
  coltitle=qftblue
}

\newtcolorbox{remarkbox}[1][]{
  enhanced,
  breakable,
  colback=white,
  colframe=qftgreen,
  boxrule=0.7pt,
  arc=1.5mm,
  title={#1},
  fonttitle=\bfseries,
  coltitle=qftgreen
}

\newtcolorbox{warningbox}[1][]{
  enhanced,
  breakable,
  colback=white,
  colframe=qftred,
  boxrule=0.7pt,
  arc=1.5mm,
  title={#1},
  fonttitle=\bfseries,
  coltitle=qftred
}

\begin{document}
\pagenumbering{gobble}

\begin{titlepage}
\centering
\vspace*{1.2in}
{\scshape\Large PHYS 662 --- Quantum Field Theory I\par}
\vspace{0.35in}
{\Huge\bfseries\color{qftblue} Lecture 3\par}
\vspace{0.15in}
{\LARGE\bfseries Quantum Sound and Phonons\par}
\vspace{0.12in}
{\Large\color{qftgreen} From Coupled Atoms to an Emergent Quantum Field\par}
\vspace{0.6in}
{\large Nima Lashkari\par}
{\large Fall 2026\par}
\vfill
\begin{minipage}{0.86\textwidth}
\small
A solid is microscopically a many-body system, but its long-wavelength
vibrations are described by a field.  Diagonalizing the harmonic lattice
reveals independent normal modes; quantizing those modes produces phonons,
the bosonic quanta of sound.  The continuum limit is a massless scalar field
theory with an emergent limiting speed, while the lattice supplies a physical
short-distance cutoff and anharmonicity supplies interactions.

\medskip
Course page: \url{https://nimalashkari.com/quantum-field-theory-2026-part-i/}
\end{minipage}
\vspace{0.5in}
\end{titlepage}

\pagenumbering{roman}
\tableofcontents
\newpage
\pagenumbering{arabic}

\section*{Roadmap}
\addcontentsline{toc}{section}{Roadmap}

Lectures 1 and 2 motivated field theory as a long-distance language and
explained why particles need not be fundamental.  This lecture develops the
first complete example.  We start with a ring of coupled atoms, find its
classical normal modes, quantize them, and then take the continuum limit.
The calculation answers five questions:
\begin{enumerate}
  \item How does a local many-particle system produce a field at long
  wavelength?
  \item Why are the collective excitations particles, and why are they
  bosons?
  \item What is the precise dispersion relation, and where does the speed of
  sound come from?
  \item How does a lattice regulate the ultraviolet continuum theory?
  \item What changes when the crystal is anharmonic and phonons interact?
\end{enumerate}

The one-dimensional chain keeps the algebra transparent.  Its longitudinal
displacement becomes a scalar field.  A realistic three-dimensional solid
has vector displacements, several polarizations, and possibly several atoms
per unit cell, but the logic is the same.

\section{Emergence and the particle--field question}

An emergent phenomenon is a robust collective behavior that is not naturally
assigned to one microscopic constituent.  Electric current in a metal, for
example, depends on the electronic many-body state, the lattice, and
scattering.  Ohm's law is not a property of an isolated atom.  In the same
way, sound is not the motion of one atom through a solid.  It is a coherent
disturbance of many atoms.

This observation changes the particle--wave question.  A photon is a quantum
of the electromagnetic field; a phonon is a quantum of the displacement
field.  A wave profile and a particle excitation are therefore not rival
microscopic pictures.  They are two descriptions of states in the same
quantum field theory:
\begin{center}
\begin{tabular}{>{\bfseries}p{0.25\textwidth} p{0.58\textwidth}}
Field language & Local amplitudes and correlation functions in space and
time. \\
Particle language & Quanta of normal modes when the spectrum contains
sufficiently long-lived excitations. \\
\end{tabular}
\end{center}

At weak coupling and low density, separated wave packets make the particle
description exceptionally useful.  When interactions remain strong, a
particle basis may cease to be useful even though local fields and their
correlators remain meaningful.

\begin{keyidea}[The example in one sentence]
The ultraviolet variables are the atomic displacements $u_j(t)$; the
infrared variable is a smooth field $u(x,t)$; and the quanta of its normal
modes are emergent particles called phonons.
\end{keyidea}

\section{A microscopic model of a solid}

\subsection{Coordinates and assumptions}

Consider $N$ identical masses $m$ arranged on a ring of circumference
\begin{equation}
  L=Na,
\end{equation}
where $a$ is the equilibrium lattice spacing.  The equilibrium position of
site $j$ is $x_j=ja$, and the actual position is
\begin{equation}
  X_j(t)=x_j+u_j(t).
  \label{eq:actualposition}
\end{equation}
The variable $u_j$ is a longitudinal displacement, not the full position.
We impose periodic boundary conditions
\begin{equation}
  u_{j+N}=u_j,
  \qquad
  p_{j+N}=p_j,
\end{equation}
and take $p_j=m\dot u_j$.

Suppose the energy is local and dominated by nearest-neighbor bonds.  If
$U(r)$ is the energy of a bond of length $r$, expand around its equilibrium
length $a$:
\begin{equation}
  U(a+\Delta u)
  =U(a)+U'(a)\Delta u
  +\frac{U''(a)}{2}(\Delta u)^2
  +\frac{U'''(a)}{3!}(\Delta u)^3+\cdots,
  \qquad
  \Delta u=u_{j+1}-u_j.
  \label{eq:bondexpansion}
\end{equation}
Mechanical equilibrium gives $U'(a)=0$.  Defining the spring constant
\begin{equation}
  K=U''(a)>0,
\end{equation}
the harmonic Hamiltonian is
\begin{equation}
  \boxed{
  H_{\mathrm{harm}}
  =\sum_{j=0}^{N-1}
  \left[
    \frac{p_j^2}{2m}
    +\frac{K}{2}(u_{j+1}-u_j)^2
  \right].
  }
  \label{eq:latticeH}
\end{equation}
The harmonic approximation is controlled when the typical relative
displacement is small enough that the cubic and higher terms in
Eq.~\eqref{eq:bondexpansion} are subleading.

\subsection{Why there is no term proportional to \texorpdfstring{$u_j$}{u j}}

A uniform translation $u_j\mapsto u_j+c$ moves the whole solid without
stretching any bond.  The Hamiltonian therefore depends on differences of
displacements, not on $u_j$ itself.  This exact symmetry forbids an onsite
potential such as $\sum_j u_j^2$ for a freely translating chain.

A substrate or pinning potential would break this symmetry and add, for
example,
\begin{equation}
  H_{\mathrm{pin}}=\frac{m\Omega_0^2}{2}\sum_j u_j^2.
\end{equation}
It would change the low-momentum dispersion from gapless to gapped.  The
absence of pinning is therefore the microscopic reason that acoustic sound
has $\omega_k\to0$ as $k\to0$.

\begin{remarkbox}[Three related facts]
Expanding about equilibrium removes the linear term in each bond extension.
Periodic boundary conditions also make
$\sum_j(u_{j+1}-u_j)=0$.  Most importantly, uniform-translation symmetry
forbids any energy cost for the zero-momentum displacement.  These are
different views of the same physical structure.
\end{remarkbox}

\section{Normal modes and the lattice dispersion relation}

\subsection{Discrete Fourier transform}

Translation invariance suggests diagonalizing the Hamiltonian in crystal
momentum.  Choose one representative of each allowed momentum in the first
Brillouin zone,
\begin{equation}
  k=\frac{2\pi n}{L},
  \qquad
  -\frac{\pi}{a}<k\leq\frac{\pi}{a}.
  \label{eq:allowedk}
\end{equation}
Define
\begin{align}
  u_k&=\frac{1}{\sqrt N}\sum_{j=0}^{N-1}
  \e^{-\ii kx_j}u_j,
  &
  u_j&=\frac{1}{\sqrt N}\sum_k \e^{\ii kx_j}u_k,
  \label{eq:DFTu}
  \\
  p_k&=\frac{1}{\sqrt N}\sum_{j=0}^{N-1}
  \e^{-\ii kx_j}p_j,
  &
  p_j&=\frac{1}{\sqrt N}\sum_k \e^{\ii kx_j}p_k.
  \label{eq:DFTp}
\end{align}
The orthogonality relation is
\begin{equation}
  \frac{1}{N}\sum_{j=0}^{N-1}
  \e^{\ii(k-k')x_j}=\delta_{k,k'}.
  \label{eq:DFTorthogonality}
\end{equation}
Because $u_j$ and $p_j$ are real classically,
\begin{equation}
  u_k^*=u_{-k},
  \qquad
  p_k^*=p_{-k}.
  \label{eq:reality}
\end{equation}
The Fourier variables have Poisson bracket
\begin{equation}
  \pb{u_k}{p_{k'}}=\delta_{k,-k'}.
  \label{eq:FourierPB}
\end{equation}
The minus sign in the Kronecker delta is a direct consequence of using the
same Fourier phase convention for $u_k$ and $p_k$.

\subsection{Diagonalizing the Hamiltonian}

The discrete derivative acts diagonally in momentum space:
\begin{equation}
  u_{j+1}-u_j
  =\frac{1}{\sqrt N}\sum_k
  \e^{\ii kx_j}(\e^{\ii ka}-1)u_k.
\end{equation}
Consequently,
\begin{equation}
  \sum_j(u_{j+1}-u_j)^2
  =\sum_k 4\sin^2\left(\frac{ka}{2}\right)u_ku_{-k}.
\end{equation}
The Hamiltonian becomes
\begin{equation}
  H_{\mathrm{harm}}
  =\sum_k\left[
    \frac{p_kp_{-k}}{2m}
    +\frac{m\omega_k^2}{2}u_ku_{-k}
  \right],
  \label{eq:Hnormal}
\end{equation}
with the exact lattice dispersion relation
\begin{equation}
  \boxed{
  \omega_k^2
  =\frac{4K}{m}\sin^2\left(\frac{ka}{2}\right),
  \qquad
  \omega_k=2\sqrt{\frac Km}
  \left|\sin\left(\frac{ka}{2}\right)\right|.
  }
  \label{eq:dispersion}
\end{equation}
Thus the coupled system is equivalent, at the harmonic level, to a
collection of decoupled normal modes.

\subsection{Equations of motion}

Hamilton's equations give
\begin{equation}
  \dot u_k=\frac{p_k}{m},
  \qquad
  \dot p_k=-m\omega_k^2u_k,
\end{equation}
and hence
\begin{equation}
  \ddot u_k+\omega_k^2u_k=0.
  \label{eq:modeEOM}
\end{equation}
Equivalently, the site variables satisfy
\begin{equation}
  m\ddot u_j
  =K(u_{j+1}-2u_j+u_{j-1}).
  \label{eq:latticeEOM}
\end{equation}
A trial wave $u_j(t)\propto\e^{\ii(kx_j-\omega t)}$ inserted into
Eq.~\eqref{eq:latticeEOM} reproduces Eq.~\eqref{eq:dispersion}.

\subsection{Long wavelengths and the speed of sound}

At momenta much smaller than the inverse lattice spacing,
$|k|a\ll1$, the exact dispersion becomes
\begin{align}
  \omega_k^2
  &=c_s^2k^2\left[1-\frac{(ka)^2}{12}
    +\mathcal O((ka)^4)\right],
  \label{eq:omegasqexpansion}
  \\
  \omega_k
  &=c_s|k|\left[1-\frac{(ka)^2}{24}
    +\mathcal O((ka)^4)\right],
  \label{eq:omegaexpansion}
\end{align}
where
\begin{equation}
  \boxed{c_s=a\sqrt{\frac Km}}
  \label{eq:soundspeed}
\end{equation}
is the speed of sound.  The linear term describes nondispersive waves.  The
higher powers of $ka$ remember the lattice and become important near the
edge of the Brillouin zone.

The maximum \emph{single-mode frequency} in this monatomic nearest-neighbor
model is
\begin{equation}
  \omega_{\max}=2\sqrt{\frac Km}.
  \label{eq:omegamax}
\end{equation}
This is not a maximum energy for the full system: arbitrarily many phonons
can occupy the available modes.

\begin{warningbox}[Dimensions and notation]
Here $K$ is a spring constant with units of energy per length squared.  With
this convention $\omega_k=2\sqrt{K/m}\,|\sin(ka/2)|$ and
$c_s=a\sqrt{K/m}$.  If one instead writes the bond energy as
$\lambda[(u_{j+1}-u_j)/a]^2/2$, then $\lambda=Ka^2$ and
$c_s=\sqrt{\lambda/m}$.  Mixing the two conventions loses a factor of $a$.
\end{warningbox}

\section{Quantization: phonons are quanta of the normal modes}

\subsection{Canonical commutators and the zero mode}

Quantization promotes the lattice variables to operators obeying
\begin{equation}
  \comm{\hat u_j}{\hat p_\ell}=\ii\hbar\delta_{j\ell},
  \qquad
  \comm{\hat u_j}{\hat u_\ell}
  =\comm{\hat p_j}{\hat p_\ell}=0.
  \label{eq:latticeCCR}
\end{equation}
In Fourier space,
\begin{equation}
  \comm{\hat u_k}{\hat p_{k'}}
  =\ii\hbar\delta_{k,-k'},
  \qquad
  \hat u_k^\dagger=\hat u_{-k},
  \qquad
  \hat p_k^\dagger=\hat p_{-k}.
  \label{eq:FourierCCR}
\end{equation}

The $k=0$ mode requires special care.  It has $\omega_0=0$ and represents a
uniform translation, not a harmonic oscillator.  If
\begin{equation}
  U_{\mathrm{cm}}=\frac1N\sum_j u_j,
  \qquad
  P=\sum_jp_j,
\end{equation}
then $u_0=\sqrt N\,U_{\mathrm{cm}}$, $p_0=P/\sqrt N$, and
\begin{equation}
  H_{\mathrm{cm}}=\frac{P^2}{2Nm}.
  \label{eq:Hcm}
\end{equation}
It is a free center-of-mass degree of freedom.  One may fix the center of
mass, restrict to $P=0$, or add weak pinning.  One must not insert
$\omega_0=0$ into oscillator formulas containing $1/\sqrt{\omega_0}$.

\subsection{Creation and annihilation operators}

For every nonzero $k$, define
\begin{align}
  \hat a_k
  &=\sqrt{\frac{m\omega_k}{2\hbar}}\,\hat u_k
  +\frac{\ii}{\sqrt{2\hbar m\omega_k}}\,\hat p_k,
  \label{eq:akdef}
  \\
  \hat a_k^\dagger
  &=\sqrt{\frac{m\omega_k}{2\hbar}}\,\hat u_{-k}
  -\frac{\ii}{\sqrt{2\hbar m\omega_k}}\,\hat p_{-k}.
  \label{eq:akdagdef}
\end{align}
Then
\begin{equation}
  \comm{\hat a_k}{\hat a_{k'}^\dagger}=\delta_{k,k'},
  \qquad
  \comm{\hat a_k}{\hat a_{k'}}=0.
  \label{eq:bosonCCR}
\end{equation}
The inverse relations make the pairing of $k$ and $-k$ explicit:
\begin{align}
  \hat u_k
  &=\sqrt{\frac{\hbar}{2m\omega_k}}
  \left(\hat a_k+\hat a_{-k}^\dagger\right),
  \label{eq:ukladder}
  \\
  \hat p_k
  &=-\ii\sqrt{\frac{\hbar m\omega_k}{2}}
  \left(\hat a_k-\hat a_{-k}^\dagger\right).
  \label{eq:pkladder}
\end{align}
This pairing is required by the Hermiticity of the real-space displacement.
It does not identify $\hat a_k$ with $\hat a_{-k}$: right- and left-moving
phonons are independent excitations.

Substitution into Eq.~\eqref{eq:Hnormal} gives
\begin{equation}
  \boxed{
  \hat H
  =\frac{\hat P^2}{2Nm}
  +\sum_{k\neq0}\hbar\omega_k
  \left(\hat a_k^\dagger\hat a_k+\frac12\right).
  }
  \label{eq:quantumH}
\end{equation}
The number operator
\begin{equation}
  \hat N_k=\hat a_k^\dagger\hat a_k
\end{equation}
counts phonons of crystal momentum $\hbar k$ in the harmonic theory.

\begin{keyidea}[The first QFT]
Each normal mode is a quantum harmonic oscillator, but the collection of all
modes is a quantum field.  The creation operator $\hat a_k^\dagger$ adds one
quantum of a collective vibration.  That quantum is a phonon.
\end{keyidea}

\section{One phonon, many phonons, and Fock space}

Choose the translational sector and define the phonon vacuum by
\begin{equation}
  \hat a_k\ket{\Omega}=0,
  \qquad k\neq0.
  \label{eq:phononvacuum}
\end{equation}
A one-phonon momentum eigenstate is
\begin{equation}
  \ket{k}=\hat a_k^\dagger\ket{\Omega},
  \qquad
  \hat H\ket{k}=(E_0+\hbar\omega_k)\ket{k}.
  \label{eq:onephonon}
\end{equation}
A localized phonon is not an additional fundamental object; it is a wave
packet of momentum modes,
\begin{equation}
  \ket{f}
  =\sum_{k\neq0}f_k\hat a_k^\dagger\ket{\Omega},
  \qquad
  \sum_{k\neq0}|f_k|^2=1.
  \label{eq:wavepacket}
\end{equation}
Because $k$ is restricted to a Brillouin zone, no one-phonon state can be
localized on scales parametrically smaller than the lattice spacing.

Multiphonon states are
\begin{equation}
  \ket{\{n_k\}}
  =\prod_{k\neq0}
  \frac{(\hat a_k^\dagger)^{n_k}}{\sqrt{n_k!}}
  \ket{\Omega}.
  \label{eq:numberstates}
\end{equation}
Since creation operators commute, exchanging two phonons leaves the state
unchanged.  Phonons are bosons.  If $\cH_1$ denotes the one-phonon Hilbert
space, the full state space is the symmetric Fock space
\begin{equation}
  \cF_s(\cH_1)
  =\bigoplus_{n=0}^{\infty}
  \operatorname{Sym}\left(\cH_1^{\otimes n}\right).
  \label{eq:Fockphonon}
\end{equation}
This does not mean that every state has an indefinite number of phonons.
Fock space contains every fixed-number sector and also permits their
superpositions.

\subsection{Microscopic atoms versus emergent phonons}

The normal-mode transformation reorganizes the same underlying system:
\begin{center}
\begin{tabular}{>{\bfseries}p{0.25\textwidth} p{0.29\textwidth} p{0.31\textwidth}}
 & Microscopic description & Infrared description \\
\toprule
Variables & Site displacements $u_j$ & Smooth displacement field $u(x)$ \\
Excitations & Atomic motion & Phonons \\
Dispersion & Set by microscopic dynamics & $E_k\simeq\hbar c_s|k|$ \\
Statistics & Not inferred from site labels & Bosonic normal-mode quanta \\
Resolution & Lattice scale $a$ & Wavelength $\lambda\gg a$ \\
\bottomrule
\end{tabular}
\end{center}

Calling the site variables ``distinguishable atoms'' is only shorthand for
the fact that equilibrium sites label the normal coordinates.  Actual atoms
of the same isotope are identical quantum particles.  Their microscopic
exchange statistics and the bosonic statistics of phonons are logically
different statements.

\section{The continuum limit: an emergent scalar field}

\subsection{From sums and differences to an integral}

For a slowly varying displacement,
\begin{equation}
  u_{j+1}-u_j
  =a\,\partial_xu(x_j)+\mathcal O(a^2),
  \qquad
  \sum_j f(x_j)\longrightarrow\frac1a\int_0^L\dd x\,f(x).
  \label{eq:continuumrules}
\end{equation}
Introduce the mass density and one-dimensional elastic coefficient
\begin{equation}
  \rho=\frac ma,
  \qquad
  C=Ka.
  \label{eq:rhoC}
\end{equation}
Then Eq.~\eqref{eq:latticeH} becomes
\begin{equation}
  H
  =\int_0^L\dd x\left[
    \frac{\rho}{2}(\partial_tu)^2
    +\frac{C}{2}(\partial_xu)^2
  \right]
  +\text{higher-derivative terms}.
  \label{eq:Hcontinuumu}
\end{equation}
The continuum canonical momentum density is
\begin{equation}
  p_u(x)=\rho\,\partial_tu(x).
\end{equation}

It is convenient to canonically normalize the field,
\begin{equation}
  \phi(x,t)=\sqrt\rho\,u(x,t),
  \qquad
  \Pi(x,t)=\partial_t\phi(x,t).
  \label{eq:phirescale}
\end{equation}
Using $c_s^2=C/\rho=Ka^2/m$, the Hamiltonian is
\begin{equation}
  \boxed{
  H_{\mathrm{IR}}
  =\frac12\int\dd x\left[
    \Pi^2+c_s^2(\partial_x\phi)^2
  \right].
  }
  \label{eq:Hscalar}
\end{equation}
The corresponding action is
\begin{equation}
  S_{\mathrm{IR}}[\phi]
  =\frac12\int\dd t\,\dd x\left[
    (\partial_t\phi)^2-c_s^2(\partial_x\phi)^2
  \right],
  \label{eq:Sscalar}
\end{equation}
and the equation of motion is the wave equation
\begin{equation}
  \boxed{
  (\partial_t^2-c_s^2\partial_x^2)\phi(x,t)=0.
  }
  \label{eq:waveequation}
\end{equation}

\subsection{Continuum quantization and the cutoff}

The equal-time canonical commutator is
\begin{equation}
  \comm{\hat\phi(x,t)}{\hat\Pi(y,t)}
  =\ii\hbar\delta(x-y).
  \label{eq:continuumCCR}
\end{equation}
With the center-of-mass mode treated separately, a continuum mode expansion
is
\begin{align}
  \hat\phi(x,t)
  &=\int_{|k|<\Lambda}\frac{\dd k}{2\pi}
  \sqrt{\frac{\hbar}{2\omega_k}}
  \left[
    \hat a(k)\e^{\ii(kx-\omega_kt)}
    +\hat a^\dagger(k)\e^{-\ii(kx-\omega_kt)}
  \right],
  \label{eq:continuumphi}
  \\
  \hat\Pi(x,t)
  &=-\ii\int_{|k|<\Lambda}\frac{\dd k}{2\pi}
  \sqrt{\frac{\hbar\omega_k}{2}}
  \left[
    \hat a(k)\e^{\ii(kx-\omega_kt)}
    -\hat a^\dagger(k)\e^{-\ii(kx-\omega_kt)}
  \right],
  \label{eq:continuumPi}
\end{align}
where
\begin{equation}
  \comm{\hat a(k)}{\hat a^\dagger(k')}
  =2\pi\delta(k-k').
  \label{eq:continuumaCCR}
\end{equation}
For the exact lattice theory, $\Lambda$ is of order $\pi/a$ and
$\omega_k$ is the sine dispersion in Eq.~\eqref{eq:dispersion}.  The
strict infrared field theory replaces it by $\omega_k=c_s|k|$ and is valid
only for $|k|a\ll1$.

At finite cutoff, Eq.~\eqref{eq:continuumCCR} contains a regulated delta
function.  The ordinary Dirac delta is recovered only after modes with
arbitrarily large momentum are formally restored.  The crystal therefore
gives a concrete ultraviolet regulator: it limits both the available momenta
and the spatial resolution.

\begin{remarkbox}[What the continuum limit means]
Taking $a\to0$ is a useful idealization, but a real solid retains a nonzero
lattice spacing.  The continuum theory is an effective theory: it accurately
predicts long-wavelength observables and organizes corrections in powers of
$ka$, while the microscopic lattice determines its coefficients and its
breakdown scale.
\end{remarkbox}

\subsection{Emergent Lorentz symmetry}

If we use $c_s$ as the limiting speed, Eq.~\eqref{eq:waveequation} has the
same form as the massless Klein--Gordon equation in $1+1$ dimensions.  Its
leading low-energy dynamics is invariant under Lorentz transformations built
with $c_s$, and its solutions split into left- and right-moving waves,
\begin{equation}
  \phi(x,t)=f(x-c_st)+g(x+c_st).
\end{equation}

This Lorentz symmetry is emergent, not exact.  The solid has a preferred rest
frame, and the corrections in Eq.~\eqref{eq:omegaexpansion} violate the
linear relativistic dispersion.  Near $|k|\sim\pi/a$, the lattice is visible
and the symmetry is lost.

\begin{warningbox}[Do not confuse the two limiting speeds]
The emergent symmetry of sound uses $c_s$, not the speed of light $c$.
Microscopically the solid is a nonrelativistic medium embedded in a
relativistic world.  Approximate Lorentz invariance can emerge in the
infrared even when it is absent from the effective microscopic lattice
model.
\end{warningbox}

\subsection{The physical field in one dimension}

Uniform shifts of $u$ are unobservable in the unpinned theory.  The local
strain
\begin{equation}
  \varepsilon(x,t)=\partial_xu(x,t)
\end{equation}
and the velocity $\partial_tu$ are directly physical.  In one spatial
dimension the massless scalar itself also has infrared subtleties associated
with its zero mode, while derivative observables are well behaved.  This is
the continuum counterpart of separating the center of mass before defining
phonon ladder operators.

\section{Anharmonicity and interacting phonons}

The harmonic Hamiltonian is exactly solvable because it is quadratic.  Real
bonds are not perfectly harmonic.  Write
\begin{equation}
  U(a+\Delta u)
  =U(a)+\frac K2(\Delta u)^2
  +\frac{g_3}{3!}(\Delta u)^3
  +\frac{g_4}{4!}(\Delta u)^4+\cdots.
  \label{eq:anharmonicpotential}
\end{equation}
If the bond potential is symmetric about its minimum, $g_3=0$; otherwise a
cubic term is generally present.  In momentum space the cubic interaction is
\begin{align}
  H_3
  ={}&\frac{g_3}{3!\sqrt N}
  \sum_{k_1,k_2,k_3}
  \prod_{r=1}^3(\e^{\ii k_ra}-1)
  u_{k_1}u_{k_2}u_{k_3}
  \delta^{(G)}_{k_1+k_2+k_3,0},
  \label{eq:H3momentum}
\end{align}
where the delta symbol enforces crystal-momentum conservation modulo a
reciprocal lattice vector
\begin{equation}
  G=\frac{2\pi\ell}{a},
  \qquad \ell\in\mathbb Z.
\end{equation}
Processes with nonzero $G$ are Umklapp processes.  Thus anharmonicity couples
normal modes, but not arbitrarily: lattice translation symmetry imposes a
selection rule.

Substituting Eqs.~\eqref{eq:ukladder}--\eqref{eq:pkladder} produces terms of
the schematic form
\begin{equation}
  \hat a\hat a\hat a,
  \qquad
  \hat a^\dagger\hat a\hat a,
  \qquad
  \hat a^\dagger\hat a^\dagger\hat a,
  \qquad
  \hat a^\dagger\hat a^\dagger\hat a^\dagger.
  \label{eq:cubicterms}
\end{equation}
These interactions can change phonon number.  For example, when energy and
crystal momentum allow it, one phonon can decay into two, or two phonons can
merge into one.  Phonon number is exactly conserved by the harmonic
Hamiltonian but is not a fundamental conservation law of the solid.

In the continuum limit the interactions are derivative interactions,
\begin{equation}
  H_{\mathrm{int}}
  \sim\int\dd x\left[
    \frac{g_3a^2}{3!}(\partial_xu)^3
    +\frac{g_4a^3}{4!}(\partial_xu)^4+\cdots
  \right].
  \label{eq:continuuminteractions}
\end{equation}
Every field appears differentiated because $u\mapsto u+c$ remains a
symmetry.  At sufficiently small momentum these higher-order interactions
are weak, which explains why long-wavelength phonons can be long lived even
when the microscopic solid is not weakly interacting.

\begin{keyidea}[Particle number is emergent]
The quadratic approximation makes each $N_k$ conserved and gives perfectly
stable phonons.  Anharmonicity breaks those separate conservation laws.
Phonons remain useful particles only when their decay widths are small
compared with their energies.
\end{keyidea}

\section{In what sense are phonons real?}

Phonons carry energy and crystal momentum, transport heat, scatter from
defects and other phonons, and appear as peaks in measurable response
functions.  A standard diagnostic is the dynamical structure factor,
schematically
\begin{equation}
  S(k,\omega)
  =\int\dd t\,\e^{\ii\omega t}
  \left\langle \hat\rho_k(t)\hat\rho_{-k}(0)\right\rangle,
  \label{eq:structurefactor}
\end{equation}
which is probed by neutron, x-ray, or light scattering.  In the ideal
harmonic theory it contains sharp contributions at the normal-mode
frequencies.  Interactions shift the frequencies and broaden the peaks.  A
narrow peak corresponds to a long-lived phonon quasiparticle.

For finite $N$, the momenta and frequencies are discrete.  In a macroscopic
solid the spacing between adjacent allowed $k$ values is $2\pi/L$ and becomes
so small that the band appears continuous.  The single-phonon band remains
bounded because the Brillouin zone is compact, even though the many-phonon
Hilbert space has no upper energy bound.

\begin{remarkbox}[Reality does not require fundamentality]
Phonons are not among the microscopic constituents of a solid, but they are
real excitations of the solid in precisely the regime where the corresponding
spectral peaks are sharp.  ``Emergent'' means regime-dependent, not fictitious.
\end{remarkbox}

\section{Beyond the one-dimensional monatomic chain}

In three dimensions each atom has a vector displacement.  For a crystal with
one atom per primitive cell, diagonalizing the dynamical matrix gives three
acoustic branches near $\bm k=0$: one longitudinal and two transverse.  If a
primitive cell contains more than one atom, additional optical branches are
possible.  Their frequencies need not vanish at $\bm k=0$.

The scalar field in Eq.~\eqref{eq:Sscalar} therefore captures one
longitudinal polarization in one dimension.  The general elastic theory uses
a displacement vector $u_i(\bm x,t)$ and a strain tensor
\begin{equation}
  u_{ij}=\frac12(\partial_i u_j+\partial_j u_i).
\end{equation}
The elastic constants determine the velocities and polarizations of the
different branches.  Quantizing every normal mode again produces phonons.

\section{Lessons for quantum field theory}

This example illustrates several general principles.
\begin{enumerate}
  \item \textbf{Locality.}  Nearest-neighbor forces produce a local
  derivative expansion in the continuum.
  \item \textbf{Symmetry.}  Discrete translation diagonalizes the harmonic
  problem; uniform-displacement symmetry protects the gapless mode; and
  crystal momentum is conserved modulo reciprocal lattice vectors.
  \item \textbf{Positivity and stability.}  $K>0$ makes the harmonic energy
  nonnegative and the normal-mode frequencies real.
  \item \textbf{Coarse-graining.}  At $|k|a\ll1$, the discrete displacement
  becomes a smooth field and microscopic details enter through a few
  coefficients such as $\rho$, $C$, and the anharmonic couplings.
  \item \textbf{Particles from fields.}  Quantizing normal modes creates a
  Fock space of indistinguishable bosonic quanta.
  \item \textbf{Effective symmetry.}  The infrared wave equation has an
  emergent Lorentz symmetry that the lattice does not possess.
  \item \textbf{Regime-dependent particles.}  The harmonic theory has exact
  phonons; the interacting theory has quasiparticles whose utility is set by
  their lifetime.
\end{enumerate}

The same logic appears throughout QFT.  A microscopic theory need not be a
field theory of the variables that are simplest in the infrared.  Conversely,
a continuum field theory should be understood together with its regime of
validity and the operators that encode corrections from shorter distances.

\section{Questions to think about}

\begin{enumerate}
  \item Which symmetry guarantees $\omega_{k=0}=0$ in an unpinned crystal?
  How does a substrate change the answer?
  \item Why does the Fourier-space canonical commutator contain
  $\delta_{k,-k'}$, while the ladder-operator commutator contains
  $\delta_{k,k'}$?
  \item What goes wrong if one defines $\hat a_0$ by dividing by
  $\sqrt{\omega_0}$?
  \item In what precise sense is the energy of a single phonon bounded, while
  the energy of the solid is not?
  \item Why can anharmonic interactions change phonon number without
  violating energy conservation?
  \item Which terms in the exact dispersion violate emergent Lorentz
  symmetry?
  \item Are microscopic atoms in a crystal really distinguishable?  Why is
  that question separate from the statistics of phonons?
  \item What observable feature distinguishes a long-lived phonon from a
  strongly damped collective excitation?
\end{enumerate}

\section*{Summary}
\addcontentsline{toc}{section}{Summary}

A periodic chain with local harmonic interactions is diagonalized by a
discrete Fourier transform.  Its exact normal-mode frequencies are
$\omega_k=2\sqrt{K/m}|\sin(ka/2)|$.  At long wavelength this becomes
$\omega_k\simeq c_s|k|$ with $c_s=a\sqrt{K/m}$, and the lattice equation
becomes the massless scalar wave equation.  Quantizing every nonzero normal
mode produces bosonic creation and annihilation operators; their quanta are
phonons, and their multiparticle state space is a symmetric Fock space.  The
uniform zero mode is a free center-of-mass coordinate and must be treated
separately.  The lattice supplies a physical ultraviolet cutoff and breaks
the Lorentz symmetry that emerges at low energy.  Anharmonic terms couple
normal modes, conserve crystal momentum modulo reciprocal lattice vectors,
and make phonon number nonconserved.  Phonons remain meaningful particles
when the resulting spectral peaks are narrow and their lifetimes are long.

\end{document}
